\section{Learned computations at retained checkpoints}
\label{note:checkpoint_computation}

This post hoc analysis makes the computations learned in main Figs.~5 and 6 visible. It uses every seed in each primary cohort and the original validation-selected states, without additional training, rate selection or selection of illustrative seeds. The two cohorts remain separate; they are not forty replicates of one experiment.

\paragraph{Single-neuron branch tuning.}
For each of the twenty opposed-target teacher--student seeds, we evaluated exact credit, the hard distal gate and broadcast at the primary Adam rate 0.03 and the state selected within 4,096 updates. Each seed's own teacher provides the reference. We varied the first feature over 65 values in $[-2,2]$ and integrated over the independent second feature using 32-point Gauss--Legendre quadrature for the uniform distribution on that interval. Context selected the left or right branch. The plotted quantity is the selected proximal branch's contribution to somatic voltage, $g_{p\to0}^{\rm den}V_p/(1+\sum_p g_{p\to0}^{\rm den})$, rather than its unweighted voltage: coupling changes can otherwise rescale an internal response without changing its somatic effect. Means and pointwise 95\% whole-seed bootstrap intervals use all twenty seeds and 10,000 resamples. The exact and gated learners recover the opposing tuning curves closely; broadcast produces weaker feature modulation (Fig.~\ref{fig:si_checkpoint_computation}A,B).

\paragraph{Population response surfaces.}
We evaluated all twenty original-bound Adam rescue seeds for resistance gating, augmentation and exact credit at their original selected rates and checkpoints. The model code is the verified execution snapshot. Each selected stream was evaluated on a uniform $25\times25$ feature grid over $[-2,2]^2$. At each grid point, sixteen fixed antithetic draws of the irrelevant streams approximate their uniform marginal distribution; the same draws are shared across rules and seeds. We multiply predictions by $(-1)^c$ to align the target sign across the four contexts, then average contexts within each seed. This makes the common target $T(z_1,z_2)=(\tanh z_1+\tanh z_2)/2+\tanh z_1\tanh z_2/4$. The plotted surfaces are seed means, not individual favorable examples. Per-context predictions are retained in Source Data.

To isolate the interaction, we subtract both one-feature marginal means and add the grand mean. With normalized trapezoidal grid weights $w$, this is
\[
I[F]_{ij}=F_{ij}-\sum_k w_kF_{ik}-\sum_k w_kF_{kj}+\sum_{k\ell}w_kw_\ell F_{k\ell}.
\]
The target interaction is $I[T]=\tanh z_1\tanh z_2/4$. Resistance gating reproduces much of the additive component but severely underestimates this interaction; augmentation and exact credit reproduce both (Fig.~\ref{fig:si_checkpoint_computation}C--J). For each seed we also report grid-weighted squared component error divided by target-component squared amplitude. Mean relative interaction error is 0.783 for resistance, $6.1\times10^{-5}$ for augmentation and $6.7\times10^{-5}$ for exact credit; corresponding additive errors are 0.00100, $1.0\times10^{-5}$ and $7.6\times10^{-6}$. These descriptive, context-averaged grid errors are not the ordinary-test NMSE and do not replace its primary comparisons. Marginalization and averaging can hide context-specific or seed-specific errors; the per-seed component values and per-context surfaces are retained. This analysis characterizes the existing selected states and does not establish a new convergence result.

Source records under \path{source_data/checkpoint_computation/} include the evaluation protocol, verified input hashes, all curves and surfaces, component errors and plotted summaries. The analysis script is \path{scripts/review_completion/checkpoint_computation.py}; the figure is rebuilt by \path{scripts/review_completion/checkpoint_figure.py}.

\clearpage
\begin{figure}[p]
\centering
\includegraphics[width=\textwidth]{supplementary/curated/checkpoint_computation.pdf}
\caption{\textbf{Checkpoint responses show preserved branch tuning and recovery of the target interaction.}
\textbf{A,B}, Selected left/right proximal branches' contributions to somatic voltage in the opposed-target single-neuron task. Feature 1 varies; feature 2 is integrated over its uniform distribution. Dashed grey, teacher; amber, broadcast; green, hard distal gate; dotted red-brown, exact credit. Exact and gated curves nearly coincide. Lines and bands are means and pointwise 95\% whole-seed bootstrap intervals over twenty seeds, using the primary rate and validation-selected states within 4,096 updates.
\textbf{C--F}, Nonlinear-population target and learned response surfaces, with a common color scale. Predictions are averaged over sixteen irrelevant-stream draws and four sign-aligned contexts within each seed, then over the twenty rescue seeds. Each rule retains its original selected Adam rate and checkpoint. Feature axes span the ordinary training range.
\textbf{G--J}, The same surfaces after subtracting both one-feature marginal means and adding the grand mean. The shared interaction scale is distinct from the response scale above. Resistance gating underestimates the interaction; augmented and exact credit recover it. Colors encode response values, not learning-rule identity.
These are post hoc descriptions of existing checkpoints, with no new training cohort or outcome-based choice of examples. Full per-seed and per-context values are in Source Data: \texttt{checkpoint\_computation/}.}
\label{fig:si_checkpoint_computation}
\end{figure}
\clearpage
